\documentclass{article}
\usepackage{amsmath,amssymb,hyperref}
\title{Transporting an interval proof across a change of causal order}
\begin{document}
\maketitle
\paragraph{Interface.}
Fix a finite set \(E\) of retained events with an old partial order \(P\), and a candidate order \(P'\) on \(E'\supseteq E\). Retained event identities and meaning are fixed. Every raw grant is a box \(\ell\subseteq X\subseteq h\) on \(E\). Old validation exported only the extrema of the satisfying old ideals, or an empty summary. A subsequent consumer tests eligibility of a cut and compatibility with any new box inside \(E'\). Cuts represent possible histories of a supplied order, not proof of present effect or availability of the complete order to the controller.

\paragraph{Exact transport.}
Uniform substitution of the old interval for the raw box is sound and complete under all such tests iff \(X'\cap E\in{\rm Id}(P)\) for every \(X'\in{\rm Id}(P')\). Sufficiency follows because the raw box and old interval agree on old ideals. This condition holds precisely when the induced order of \(P'\) on \(E\) contains \(P\): containment implies closure under old predecessors, and if \(a\le b\) was lost, the new principal ideal at \(b\) omits \(a\). For necessity suppose that a projected new cut \(Y\) is not an old ideal. The two requirements to include \(Y\), and to include its old downward closure, coincide on all old ideals, so have the same old summary. The new cut projecting to \(Y\) distinguishes them. Pinning it by a singleton new box shows that no method given just the old interval can answer all these queries. Necessity is uniform over all old boxes and new singleton tests; a restricted consumer can permit a weaker criterion.

\paragraph{Constructive reuse.}
For an old nonempty interval with endpoints \(L,U\) and a new raw box \(l',u'\), require \(L\cup l'\) and exclude \((E\setminus U)\cup(E'\setminus u')\). Use new edges to close the required events downward and excluded events upward. An intersection proves infeasibility; otherwise the required closure is the least witness, and the complement of the excluded closure is the greatest. Graph traversal suffices for the supplied dependency graph. Edge discovery and authentication are separate costs; a direct reader with the same verified summary and edges ties.

\paragraph{Projection and residual.}
The image of \([A,B]\) on \(E\) is exactly \([A\cap E,B\cap E]\) in the induced order of \(P'\), not necessarily in the old ambient order. For any induced ideal \(Y\) between those projected bounds, close \(Y\cup A\) downward in \(P'\). It stays in \(B\); old predecessors of \(Y\) are already in \(Y\), so its restriction is exactly \(Y\). A fiber is also an interval, since fixing \(Y\) adds the requirements \(Y\) included and \(E\setminus Y\) excluded. This preserves historical feasibility but says nothing about independently joining projections onto multiple interfaces that share hidden choices.

\paragraph{Checked separation and prior.}
A gate commit depends on a new appointment; the old license requires omission of withdrawal. Adding withdrawal as an ancestor of appointment destroys the common witness although each rule remains separately satisfiable. In a different replacement, removing the dependency distinguishes requiring just appointment from explicitly requiring both events, though their old exports agreed exactly. \href{https://users.ece.utexas.edu/~garg/dist/dcs01-slice.pdf}{Computation slicing} provides the inherited lattice of consistent cuts and precise slices of general regular predicates. This scope is reuse of a smaller interval export across changed reachability, not an improved general slicer. The graph is authored and exhaustively checked in the finite tests; no field, independent-transfer, or contract acceptance claim follows.
\end{document}
